%%=============================================================================
%% chapters/chapter2.tex — Related Work (final-defense version)
%%=============================================================================
\chapter{Related Work}
\label{ch:related}

\section{Autonomous Exploration and Next-Best-View Planning}
\label{sec:rw-exploration}

Simultaneous localization and mapping (SLAM) estimates a robot's pose while
reconstructing the geometry of its environment, progressing from filter-based
formulations such as GMapping~\cite{grisetti2007} to graph-based methods such as
Cartographer~\cite{hess2016} that enforce global consistency through loop
closure. Exploration governs how a robot expands its map and reduces
uncertainty, and has been studied principally under two
paradigms~\cite{placed2023,lluvia2021}. Frontier-based methods drive the robot
toward the boundary between known and unknown
space~\cite{yamauchi1997,umari2017}, prioritizing frontiers by expected
information gain and travel cost. Next-best-view (NBV) methods instead evaluate
candidate viewpoints directly by an information-gain or coverage
criterion~\cite{connolly1985,bircher2016}; information-theoretic formulations
make this gain explicit by quantifying the expected reduction in map
uncertainty~\cite{stachniss2005}, and the idea extends to the coordination of
multiple robots~\cite{burgard2005}. Both paradigms have been refined
extensively, for example through improved frontier extraction and cost
shaping~\cite{he2025,quin2021}, incremental frontier structures with
hierarchical planning~\cite{zhou2021}, and sampling-based online informative
path planning~\cite{schmid2020}. Most of these methods, however, assume a
sensor that acquires continuously while the platform moves and registers
incrementally online, so the dominant cost is travel and the act of observing
is effectively free along the path. Under that assumption, the question of
whether to pay a large fixed cost to stop and observe at a particular pose does
not arise. The present work departs from this assumption by introducing a
stop-and-scan exploration strategy whose scan/skip decision is made on the live
map, detailed in Chapter~\ref{ch:mapping}.

\section{Coverage Models, Visibility, and Stop-and-Scan Mapping}
\label{sec:rw-coverage}

The notion that a viewpoint observes only what is in its line of sight is
classical in computational geometry, where the visibility polygon defines the
region directly visible from a point, and in the art-gallery problem, which
asks for a minimal set of guards that jointly see an entire
domain~\cite{gonzalez2002,orourke1987,urrutia2000}. Sensor-coverage and
view-planning formulations build on this idea to place viewpoints that maximize
observed surface, and have been surveyed for automated 3D reconstruction and
inspection~\cite{scott2003}, for sensor planning in computer
vision~\cite{tarabanis1995}, and for view planning in robot active
vision~\cite{zeng2020}. Coverage path planning, which instead seeks a path that
sweeps a sensor footprint over an entire area, has likewise been surveyed
extensively~\cite{choset2001,galceran2013}, including for multi-robot model
reconstruction and mapping~\cite{almadhoun2019}. In robotic exploration
practice, however, the per-pose coverage used to drive online decisions is
frequently approximated by a range disk or a sensor cone without explicit
occlusion reasoning, because such approximations are cheap to evaluate at every
candidate. This approximation is benign when observation is continuous and
cheap, but it becomes consequential in the stop-and-scan regime, where a single
skipped occluded region corresponds to a missing multi-minute scan and an
unobserved part of the map.

Survey-grade terrestrial scanning, as practiced with tripod-class instruments,
departs from the continuous-acquisition assumption: each measurement requires a
stationary dwell of minutes, and registration is performed
offline~\cite{stopandgo-review}. Prior automated TLS and stop-and-go mapping
efforts focus largely on the system integration of driving a scanner between
viewpoints and on offline registration of the collected
scans~\cite{stopandgo-registration}, while scan-planning studies in
construction optimize viewpoint placement to maximize the visible area of
target structures~\cite{tls-scanplanning}. Some of these planning methods are
explicitly occlusion- or visibility-aware: Mozaffar and
Varshosaz~\cite{mozaffar2016occlusion} optimize scanner placement to reduce
occlusions, and Prieto et al.~\cite{prieto2017nbs} drive a next-best-scan
policy with a line-of-sight visibility analysis. Recent formulations cast
indoor scan planning as combinatorial optimization over a known floor plan:
Dehbi et al.~\cite{dehbi2021} compute minimal TLS viewpoint sets under
visibility and registrability constraints, and Knechtel et
al.~\cite{knechtel2025} jointly optimize stop-and-go standpoints and the route
between them with enforced network redundancy. Crucially, these methods operate
offline on a prior model of the environment: a building information model
(BIM), a floor plan, or an initial coarse scan. The decision of which
candidate poses are worth a stationary scan during online exploration of an
unknown space is still typically handled by fixed spacing or by a proximity
(disk) test.

The two bodies of work above leave a specific gap. NBV and frontier
exploration operate online in unknown environments but assume continuous,
low-cost acquisition. Automated TLS and scan-planning methods reason about
expensive stationary acquisition but are typically offline and assume a known
environment model. The regime addressed here is the intersection of the two:
online exploration of an unknown environment in which each acquisition carries
a large stationary cost, where the scan/skip decision must be made on the live
map. In that regime the coverage model is the decisive component, and this
dissertation argues that it must be occlusion-aware rather than proximity-based.
A broader class of sensors shares this workflow in that they require stationary
acquisition or a dwell time at each viewpoint, including structured-light 3D
scanners, thermal panorama units, ground-penetrating radar, gas and spectral
sensors, and non-destructive testing probes~\cite{gas-coverage,robotic-ndt};
the line-of-sight model developed here applies to the optical, non-penetrating
members of this family.

\section{3D Semantic Mapping, Scene Graphs, and the TOSM Lineage}
\label{sec:rw-semantic}

Object-level semantic mapping has progressed from object-oriented SLAM and
dense semantic fusion~\cite{salas2013slampp,mccormac2017semanticfusion} to
metric--semantic pipelines such as Kimera~\cite{rosinol2020kimera} and
real-time hierarchical scene graphs~\cite{hughes2022hydra}. 3D scene graphs
organize metric reconstructions into object-level nodes with semantic
attributes and inter-object relations~\cite{armeni20193dsg}, and recent
open-vocabulary variants attach language-queryable features to nodes and
edges~\cite{koch2024open3dsg,gu2024conceptgraphs,werby2024}, enabling
large-language-model (LLM)-driven mobile manipulation over the
graph~\cite{honerkamp2024momallm}.
Neural implicit language-embedding fields such as LERF~\cite{kerr2023lerf}
provide language-queryable representations but are neither symbolic nor
structured, making them difficult to persist as a database that downstream
services can query and verify directly. On the knowledge side, semantic maps
have been coupled to task planning since Galindo et
al.~\cite{galindo2008task}, grounded in spatial databases~\cite{deeken2018semap},
and served from ontology-backed infrastructures such as
KnowRob~\cite{tenorth2013knowrob}; see~\cite{manzoor2021survey} for an
application-oriented survey of ontology-based robot knowledge. These systems
target flexible construction; verification of what enters the graph and
maintenance after construction receive less attention.

Within the Triplet Ontological Semantic Model (TOSM) lineage that this work
extends, the original semantic modeling framework~\cite{joo2020tosm} defined
the symbolic/explicit/implicit triplet---a Symbolic layer of identifiers and
category labels, an Explicit layer of sensor-measurable pose, size, and shape,
and an Implicit layer of knowledge-derived information such as movability,
key-object indication, and spatial relations. TOSMNav~\cite{joo2022fsom}
scaled it to a navigation framework in which an OWL/SWRL knowledge base feeds
hierarchical PDDL planning, with per-robot on-demand databases,
kinematics-aware traversability rules, and re-planning with map maintenance,
validated across multi-floor and outdoor sites with four robot types. Its
successor DK-SMF~\cite{joo2025dksmf,joo2025phd} automates object and place
modeling from RGB-D exploration, combining zero-shot detection and visual
question answering with domain-knowledge filtering and LLM-derived place
semantics over SLIC-segmented regions. A recent Robot-as-a-Service
(RaaS)-oriented semantic task management framework~\cite{kwon2025}
normalizes robot, space, and task
information into TOSM and allocates tasks by cost-aware reasoning, and a
multilayer topological semantic map framework~\cite{kwon2024} produces grid,
object, segmentation, and inference layers from an autonomously explored grid
map. The navigation side of the lineage assumes a trustworthy semantic
database; the framework of this dissertation supplies that database---a
verified, incrementally maintained environment model grounded in survey-grade
3D point clouds---and differs from DK-SMF on four axes, summarized in
Table~\ref{tab:dksmf}: a survey-grade TLS backbone instead of RGB-D streams;
multi-view VLM verification with escalation instead of domain-knowledge
hallucination suppression; identity-preserving incremental upserts instead of
build-once modeling; and a place layer derived from verified objects only, so
that the confidence gate propagates across layers.

\begin{table}[htbp]
  \centering
  \caption{Positioning against DK-SMF~\cite{joo2025dksmf}. The DK-SMF column
  paraphrases that paper's published description.}
  \label{tab:dksmf}
  \begin{tabularx}{\textwidth}{l X X}
    \toprule
    \textbf{Axis} & \textbf{DK-SMF} & \textbf{This work} \\
    \midrule
    Sensing & RGB-D exploration & registered TLS (BLK360) \\
    Label trust & domain-knowledge filtering & multi-view verification + escalation \\
    Lifecycle & build-once & mint-once IDs, incremental upsert \\
    Place semantics & SLIC + LLM over detections & SLIC + ring codes over verified objects \\
    \bottomrule
  \end{tabularx}
\end{table}

\section{3D Point Cloud Segmentation and Open-Vocabulary Recognition}
\label{sec:rw-segmentation}

Early efforts to attach semantics to geometric maps integrated keypoint or
convolutional detectors into SLAM pipelines, and dense systems such as
SemanticFusion~\cite{mccormac2017semanticfusion} and
Kimera~\cite{rosinol2020kimera} combined 3D reconstruction with semantic
segmentation. Most of these are limited to a fixed, pre-trained class set.
Learned 3D instance segmentation built on point-set
backbones~\cite{qi2017pointnet,qi2017pointnetpp}---Mask3D~\cite{schult2023mask3d},
SPFormer~\cite{sun2023spformer}, Point Transformer~V3~\cite{wu2024ptv3}---and
promptable point segmentation (Point-SAM~\cite{zhou2024pointsam}) dominate
benchmark suites but inherit their training distributions; unsupervised
alternatives such as Part2Object~\cite{shi2024part2object} relax the label
requirement but still target curated indoor scenes. Open-vocabulary approaches
mitigate the closed-set limitation: point--text embedding models such as
Uni3D~\cite{zhou2024uni3d} and PointCLIP~V2~\cite{zhu2023pointclipv2} transfer
web-scale vision--language supervision~\cite{radford2021clip} to point clouds,
and scene-level systems attach open-vocabulary features to 3D points and
instances~\cite{peng2023openscene,takmaz2023openmask3d,nguyen2024open3dis},
while Grounded SAM~\cite{ren2024} couples a grounding detector with a
promptable segmenter to segment language-specified objects. Many of these
require posed multi-view RGB input or depth-registered 2D masks; for a
survey-grade scanner that already produces a colored point cloud directly,
synthesizing such imagery introduces an artificial dependency and risks
discarding the metric accuracy of the original laser data.

On industrial TLS content, this dissertation observes the transfer breaking
down: ScanNet-pretrained segmenters hallucinate furniture on industrial
equipment, a closed-set semantic head labels a wall-removed scene as
\texttt{wall}, PointCLIP~V2 collapses to a single class, and Uni3D's accuracy
\emph{drops} as the vocabulary is narrowed toward the true class list
(Section~\ref{sec:exp-seg}). The framework therefore uses a deliberately simple
geometric decomposition---seeded plane
removal~\cite{fischler1981ransac,schnabel2007efficient} plus fixed-parameter
density clustering~\cite{ester1996dbscan}---positioned not as a novel
segmentation algorithm but as a reproducible decomposition whose determinism
underwrites verification carry-over and incremental diffing, and treats
open-vocabulary output as a provisional prior to be verified rather than as a
label source.

\section{Multi-View VLM Verification and Knowledge-Grounded Querying}
\label{sec:rw-vlm}

Aggregating rendered views is a classic route to 3D shape recognition
(MVCNN~\cite{su2015mvcnn}), and using vision--language models (VLMs) to label
rendered 3D content is increasingly common in open-set scene
graphs~\cite{gu2024conceptgraphs} and language-driven
maps~\cite{huang2023vlmaps}. The ICTC~2026 study of this
dissertation~\cite{kim2026ictc} established multi-view forced-schema VLM
classification of TLS object clusters with confidence-weighted late fusion,
and showed that early fusion---all views in one prompt---degrades accuracy on
sparse point renders. This dissertation extends that mechanism into a
verification system: automatic escalation of split votes to a wider view set,
explicit retention of unresolved objects as unverified, tier-precision
analysis including unanimous-but-wrong failure cases, and integration of the
resulting status into gating, update, and place layers.

Grounding LLM answers in retrieved structure reduces but does not eliminate
unsupported claims~\cite{lewis2020rag,ji2023hallucination}, and
instruction-based guardrails are known to be brittle. The knowledge-graph
lineage to which the author contributed~\cite{kim2025isis} observed
hallucinated environment facts when querying ungrounded LLMs about mapped
spaces. This dissertation contributes a structural mitigation: rather than
instructing the model to respect verification status, the query-time view
demotes unverified types to \texttt{unknown} and quarantines the failed label in
a schema-documented field, an intervention in the data contract whose effect is
isolated in an ablation (Section~\ref{sec:exp-ablation}).

\section{Incremental Semantic-Map Maintenance}
\label{sec:rw-update}

Change detection in point clouds and lifelong mapping address geometric
currency~\cite{fehr2017tsdf}, and object-permanence reasoning tracks instances
through rearrangement in embodied settings~\cite{wald2019rescan}. Recent online
open-vocabulary systems bring incrementality to instance-semantic
maps~\cite{deng2026ovimap,zhai2026onlinepg,zhang2025funcgraph}, but what they
maintain is a geometric or embedding-level instance representation, not
verified semantic evidence. In the TOSM lineage, TOSMNav~\cite{joo2022fsom}
maintains its map online, updating object positions from detections and pruning
connectivity when a passage fails, which complements rather than covers the
present concern. What is rare is maintenance of a semantic knowledge graph
whose nodes carry verified labels, revision provenance, and stable robot-facing
identities across scan epochs: scan-once, build-once remains the default in
TLS-derived modeling, whose dominant line reconstructs as-built BIM or
parametric building models from static
scans~\cite{tang2010asbuilt,ochmann2016parametric,scan2bim}. The mint-once
identity policy with two-pass matching and confidence-gated node-level upserts
of Chapter~\ref{ch:semantic} targets this gap.

\section{Digital Twins, Scan-to-Simulation, and RaaS-Oriented Integration}
\label{sec:rw-digitaltwin}

Digital twins---virtual replicas of physical assets that stay synchronized
with their real counterparts---are increasingly adopted in manufacturing and
logistics for commissioning, monitoring, and operator training of mobile robot
fleets~\cite{tao2019}. A prerequisite for any robot-level twin is a simulation
environment that reproduces the geometry of the real workspace with sufficient
fidelity. In practice such environments are still largely modeled by hand.
Automatic reconstruction of as-built models from point clouds (scan-to-BIM) has
been studied extensively in the construction domain~\cite{scan2bim}, but its
output (IFC/BIM) is not directly consumable by robot simulators, and
game-engine reconstruction pipelines such as Gibson~\cite{gibson} produce
visually convincing meshes that seldom guarantee metric consistency.
Gazebo~\cite{koenig2004gazebo} has long been the default simulator in the ROS
ecosystem; NVIDIA Isaac Sim~\cite{isaacsim} adds GPU-accelerated physics,
OpenUSD scene composition, and photorealistic rendering, and has been adopted
for robot learning~\cite{orbit} and sim-to-real research~\cite{sim2real}.
Existing Isaac Sim workflows, however, assume the environment is authored
beforehand, and most reported robot digital twins mirror robot state into a
manually built scene. Table~\ref{tab:routes} contrasts the main routes from a
captured indoor space to a robot-usable simulation environment against the
pipeline of Chapter~\ref{ch:validation}.

\begin{table}[htbp]
  \centering
  \caption{Routes from a captured indoor space to a simulation environment
  (characteristics as reported by each line of work).}
  \label{tab:routes}
  \begin{tabularx}{\textwidth}{l l l X}
    \toprule
    \textbf{Approach} & \textbf{Automation} & \textbf{Turnaround} & \textbf{Metric validation} \\
    \midrule
    Manual modeling~\cite{tao2019} & manual & days per room & rarely reported \\
    Scan-to-BIM~\cite{scan2bim} & semi-automatic & minutes--hours & element-level \\
    Mesh real-to-sim~\cite{gibson} & offline pipeline & not reported & visual focus \\
    This work (Chapter~\ref{ch:validation}) & automatic & 78.5\,s & $\le$2\%; 5.5\,cm end to end \\
    \bottomrule
  \end{tabularx}
\end{table}

Prior digital-twin-based robotics studies have largely assumed RGB-D streams
or BIM inputs and have used the environmental representation mainly for
navigation. The RaaS-oriented control architecture provides a higher-level
frame for coordinating heterogeneous robots over a shared
representation~\cite{kwon2025}. In contrast to prior work, the present
dissertation supplies both the semantic representation and the digital twin
automatically from one survey-grade 3D point cloud, positions the verified
knowledge graph as the single semantic store from which the console, the
mission mediator, the twin, and the physical robot are all derived, and
validates that chain on hardware.

\section{Map Representations for Robot Tasking}
\label{sec:rw-representations}

Metric-topological maps, which combine a geometric grid map with a topological
graph, have been a core representation for indoor navigation since
Thrun~\cite{thrun1998}; free-space partitioning has traditionally used
sensor-based planning and Voronoi diagrams~\cite{choset1996}, and grid-map
room segmentation has been surveyed and made robust to
clutter~\cite{bormann2016room,luperto2022robust}. Queryable semantic
topological maps such as QueSTMaps~\cite{mehan2024} propose hierarchical,
queryable representations for 3D scene understanding, but presuppose 2D grid,
RGB-D, or neural representations as input. Table~\ref{tab:capability}
contrasts the representations a deployment can choose from: a 2D occupancy map
supports navigation but neither queryability nor semantics; a raw 3D scan adds
geometry but no task language; and a verified semantic knowledge graph adds
queryability, per-fact confidence, and, with this work, a maintenance path and a
simulation twin. The table is analytical; its query and update columns are
substantiated in Chapter~\ref{ch:experiments}.

\begin{table}[htbp]
  \centering
  \caption{Capability comparison of environment representations. The 2D-map
  and 3D-scan columns state definitional
  properties~\cite{macenski2020marathon,tang2010asbuilt}; each ``this work''
  cell names the section that substantiates it.}
  \label{tab:capability}
  \begin{tabularx}{\textwidth}{l c c X}
    \toprule
     & \textbf{2D map} & \textbf{3D scan} & \textbf{3D semantic KG (this work)} \\
    \midrule
    Metric navigation & \checkmark & \checkmark & \checkmark (derived; Section~\ref{sec:exp-missions}) \\
    Object geometry & -- & \checkmark & \checkmark (explicit models; Section~\ref{sec:backbone}) \\
    Task-language query & -- & -- & \checkmark (gated; Section~\ref{sec:exp-query}) \\
    Per-fact confidence & -- & -- & \checkmark (tiers; Section~\ref{sec:exp-verify}) \\
    Change maintenance & re-map & re-scan & node-level upsert (Section~\ref{sec:exp-update}) \\
    Place semantics & -- & -- & \checkmark (ring codes; Section~\ref{sec:place}) \\
    Simulation twin & -- & -- & \checkmark (scan-to-sim; Section~\ref{sec:exp-twin}) \\
    \bottomrule
  \end{tabularx}
\end{table}
